\section{Infrastrukturbedingte Herausforderungen bei der LLM-Anbindung}

Für die prototypische Implementierung war zunächst vorgesehen, ein durch die Hochschule bereitgestelltes und selbst gehostetes Sprachmodell über eine dokumentierte Ollama-Schnittstelle anzubinden. Der Zugriff sollte aus dem Hochschulnetz beziehungsweise über eine aktive VPN-Verbindung erfolgen. Als API-Endpunkte waren insbesondere die Ollama-Routen \texttt{/api/chat}, \texttt{/api/generate} und \texttt{/api/tags} vorgesehen. Der bereitgestellte Python-Wrapper verwendet ebenfalls diese Endpunkte und setzt das native Ollama-Request- und Response-Format voraus [11].

Bei der technischen Erprobung zeigte sich jedoch eine Abweichung zwischen der dokumentierten Schnittstelle und der tatsächlich bereitgestellten Infrastruktur. Der Aufruf der Basis-URL lieferte die OpenAPI-Dokumentation eines LiteLLM-Proxys. Die dokumentierten Ollama-Endpunkte \texttt{/api/tags} und \texttt{/api/chat} antworteten jeweils mit dem HTTP-Statuscode 404 und der Meldung \texttt{\{"detail":"Not Found"\}}. Dies weist darauf hin, dass die betreffenden Routen auf der aktuell erreichbaren Serverinstanz nicht registriert waren. Demgegenüber war der OpenAI-kompatible LiteLLM-Endpunkt \texttt{/v1/chat/completions} grundsätzlich erreichbar, beantwortete nicht authentifizierte Anfragen jedoch mit dem HTTP-Statuscode 401 und dem Hinweis auf einen fehlenden API-Schlüssel.

Da die HTTPS-Verbindung erfolgreich aufgebaut wurde und der Server strukturierte HTTP-Antworten lieferte, konnten grundlegende Probleme der Netzwerkverbindung, der DNS-Auflösung oder der TLS-Kommunikation weitgehend ausgeschlossen werden. Auch die aktive VPN-Verbindung stellte nicht die unmittelbare Fehlerursache dar. Die Testergebnisse deuten vielmehr auf eine Inkonsistenz zwischen der dokumentierten Ollama-API und der aktuell betriebenen LiteLLM-Proxy-Konfiguration hin. Ohne einen gültigen LiteLLM-Schlüssel oder eine alternative interne Ollama-Adresse war deshalb zum Zeitpunkt der Erprobung kein Zugriff auf die gehosteten Modelle möglich.

Die festgestellte Abweichung betrifft nicht die eigentliche Tutorlogik. Der Prompt-Builder erzeugt eine standardisierte Folge aus System-, Benutzer- und Assistenznachrichten, die grundsätzlich sowohl von nativen Ollama-Schnittstellen als auch von OpenAI-kompatiblen APIs verarbeitet werden kann [8]. Ebenso arbeiten die lokale Aufgabenverwaltung, die generische Hint-Policy sowie die SQLite-basierte Chat-Historie unabhängig vom verwendeten Modellbackend [2][4][5]. Backendabhängig sind dagegen insbesondere der API-Endpunkt, die Authentifizierung, einzelne Request-Parameter und die Struktur der Antwort.

Native Ollama-Anfragen verwenden beispielsweise den Endpunkt \texttt{/api/chat}, Ollama-spezifische Optionen wie \texttt{num\_predict} sowie eine Antwort im Feld \texttt{message.content}. Der LiteLLM-Proxy stellt dagegen den OpenAI-kompatiblen Endpunkt \texttt{/v1/chat/completions} bereit, verwendet Parameter wie \texttt{max\_tokens} und liefert den generierten Inhalt typischerweise unter \texttt{choices[0].message.content}. Zusätzlich kann ein Autorisierungsheader in der Form \texttt{Authorization: Bearer <API-Key>} erforderlich sein. Der ursprünglich implementierte Client ist auf das native Ollama-Format ausgerichtet [7].

Aus dieser Problematik folgt als Architekturentscheidung, die Tutor-Anwendung nicht dauerhaft an ein einzelnes LLM-Backend zu koppeln. Stattdessen sollte eine abstrakte Clientschicht eingeführt werden, welche die Tutorlogik von den technischen Besonderheiten des jeweiligen Inferenzdienstes entkoppelt. Eine mögliche Struktur umfasst separate Adapter für native Ollama- und OpenAI-kompatible LiteLLM-Schnittstellen. Die Auswahl des Backends, der Basis-URL, des Modells und gegebenenfalls des API-Schlüssels kann über Umgebungsvariablen erfolgen. Auf diese Weise bleiben Prompt-Erzeugung, Chat-Historie und Tutorsteuerung unverändert, während lediglich die Kommunikationsschicht ausgetauscht wird.

\[
\begin{aligned}
\text{Tutor-Anwendung}
&\rightarrow \text{abstrakte LLM-Client-Schnittstelle} \\
&\rightarrow
\begin{cases}
\text{Ollama-Adapter} \\
\text{LiteLLM-Adapter}
\end{cases}
\rightarrow \text{Sprachmodell}
\end{aligned}
\]

Für die Reproduzierbarkeit der Evaluation müssen daher nicht nur der Modellname, sondern auch der verwendete API-Typ, der öffentliche Modellalias, die Inferenzparameter und die Version des Proxy- beziehungsweise Inferenzdienstes dokumentiert werden. Das vorgesehene Modell \texttt{qwen3.6:27b} ist in der bereitgestellten Modellübersicht mit 27,8 Milliarden Parametern, einer Q4\_K\_M-Quantisierung und einem Kontextfenster von 262.144 Tokens aufgeführt [1]. Die bloße Verfügbarkeit in einer Ollama-Modellübersicht gewährleistet jedoch nicht, dass derselbe Name unverändert als öffentlicher Modellalias eines LiteLLM-Proxys verwendet werden kann.

Bis zur Klärung der Hochschulschnittstelle kann für Entwicklung und automatisierte Tests eine lokale Ollama-Instanz eingesetzt werden. Dabei ist zwischen funktionalen Softwaretests und der späteren Modellevaluation zu unterscheiden. Ein kleineres lokales Modell kann die korrekte Verarbeitung von Requests, die Chat-Historie, die Hilfestufen und die Fehlerbehandlung prüfen. Aussagen über die didaktische Qualität des vorgesehenen Hochschulmodells dürfen daraus jedoch nicht unmittelbar abgeleitet werden.

API-Schlüssel dürfen nicht im Quellcode oder Repository gespeichert werden. Für eine mögliche LiteLLM-Anbindung sind stattdessen Umgebungsvariablen, lokale Konfigurationsdateien außerhalb der Versionsverwaltung oder eine Secret-Verwaltung zu verwenden. Die TLS-Zertifikatsprüfung soll aktiviert bleiben. Darüber hinaus sollten Fehlermeldungen der externen Infrastruktur serverseitig protokolliert, gegenüber Studierenden jedoch auf eine allgemeine und datenschutzgerechte Meldung reduziert werden.

Die Infrastrukturproblematik verdeutlicht insgesamt die Bedeutung eines stabilen und dokumentierten API-Vertrags für reproduzierbare KI-Anwendungen. Für die vorliegende Arbeit stellt sie zugleich eine praktische Begründung für die modulare Trennung von Tutorlogik, Prompt-Erzeugung, Sitzungsverwaltung und Modellkommunikation dar. Durch diese Entkopplung kann der Prototyp unabhängig von kurzfristigen Änderungen der bereitgestellten Inferenzinfrastruktur entwickelt und getestet werden.